\subsection{Some lemmas}

Lemma 1. --- Let A be a Noetherian semi-local ring and B a finite A-algebra. Then the ring B is semi-local and Noetherian; let \(\mathfrak{m}_1, \ldots, \mathfrak{m}_n\) be its maximal ideals.

The canonical homomorphism from B to \(\prod_{i=1}^{n} \hat{\mathbf{B}}_{\mathfrak{m}_i}\) extends to an isomorphism from \(\hat{\mathbf{A}} \otimes_{\mathbf{A}} \mathbf{B}\) onto \(\prod_{i=1}^{n} \hat{\mathbf{B}}_{\mathfrak{m}_i}\) .

By IV, § 2, no. 5, cor. 3 to prop. 9, the ring B is semi-local and \(\mathfrak{m}_{A}B\) is an ideal of definition for it. By III, § 3, no. 4, Th. 3, (ii), the ring \(\hat{\mathbf{A}} \otimes_{\mathbf{A}} \mathbf{B}\) is the completion of B for the topology defined by its radical; one then applies III, § 2, no. 13, corollary to prop. 19.

Lemma 2. — Let A be a Noetherian ring and M an A-module. The canonical map from M to the product \(\prod_{\mathfrak{p} \in \operatorname{Ass}(M)} M_{\mathfrak{p}}\) is injective.

Indeed, let \(m\) be a nonzero element of M; then Ann( \(m\) ) is contained in a prime ideal \(\mathfrak{p}\) of A associated to M (IV, § 1, no. 1, prop. 2), and the image of \(m\) in \(M_{\mathfrak{p}}\) is nonzero (II, § 2, no. 2, prop. 4).

Lemma 3. — Let A be a Noetherian ring, x an element of A, M a finitely generated A-module, and \(\mathfrak{p}\) a prime ideal of A associated to M. Suppose that the homothety \(x_{\mathbf{M}}\) is injective. Let \(\mathfrak{q}\) be a prime ideal of A, minimal among those containing \(\mathfrak{p} + xA\). Then \(\mathfrak{q}\) is associated to the A-module M/xM.

Let N denote the submodule of M consisting of the elements m such that \(\mathfrak{p}m = 0\). We have N ∩ xM = xN; indeed, if an element m of M is such that \(\mathfrak{p}xm = 0\), then \(\mathfrak{p}m = 0\) since \(x_{M}\) is injective, hence \(m \in N\) . Consequently, the A-module N/xN is isomorphic to the submodule (N + xM)/xM of M/xM, and it suffices to show that \(\mathfrak{q}\) is associated to N/xN. Since \(\mathfrak{p}\) is associated to M, there exists an element m of M such that \(\mathfrak{p} = \operatorname{Ann}(m)\); we have \(m \in N\) , whence \(\mathfrak{p} = \operatorname{Ann}(N)\) and consequently \(\operatorname{Supp}(N/xN) = V(\mathfrak{p} + xA)\) by II, § 4, no. 4, corollary to prop. 18; therefore, \(\mathfrak{q}\) is associated to N/xN (IV, § 1, no. 4, Th. 2).

Lemma 4. — Let A be a discrete valuation ring, B a Noetherian local ring, and \(\rho: \mathbf{A} \to \mathbf{B}\) a local and flat homomorphism. If the ring \(\kappa_{\mathbf{A}} \otimes_{\mathbf{A}} \mathbf{B}\) is reduced, then B is reduced.

Suppose there exists a nonzero nilpotent element \(x\) of B, and let \(\pi\) be a uniformizer of A. Since we have \(\pi \mathbf{B} \subset \mathfrak{m}_{\mathbf{B}}\) , the ring B is separated for the \(\pi \mathbf{B}\) -adic topology. There therefore exists \(n \in \mathbf{N}\) and \(y \in \mathbf{B}\) with \(x = \pi^n y\) and \(y \notin \pi \mathbf{B}\) . Since B is flat over A, multiplication by \(\pi\) is injective in B. The class of \(y\) in \(\mathbf{B}/\pi \mathbf{B}\) is therefore a nonzero nilpotent element, which contradicts the hypothesis.
% semantic-fragments: legacy-input-seam
\relax{}
